%% 04c_results_surface.tex -- Owner group: G5. Rewritten for the camera-ready (2026-10-01).
%% The fitted forms (eq:models), tab:scope, tab:exponents, tab:robust, fig:diag and the serving
%% corollary live in 96_appendix_surface.tex (App. G). No float in this subsection.
\subsection{What the registered fit identifies}
\label{sec:scope}
Before training, we registered four hypotheses about a power-law fit of error over width $w$,
depth $d$ and refinement steps $T$ (App.~\ref{app:prereg}), two of them declared primary
(H-D2 and H-D3). H-D2 passes as declared; H-D3 and H-D4 fail their registered criteria; H-D1
cannot be decided, because the width exponent is not identified. App.~\ref{app:surface} gives the two fitted forms, the
exponents and the fit diagnostics. None of this affects the share in \S\ref{sec:gapclosed},
which uses no fitted model.

\paragraph{Steps and the two metrics (H-D2).}
The first registered primary test passes: the step exponent is larger for WER than for SIM-o, by
$\Delta\tau = \tau_{\mathrm{WER}} - \tau_{\mathrm{SIM}} = \Ndtau$ (CI \Ndtauci). Its sign
depends on how the error is scaled. When we refit the same measurements after a monotone
change of the error, such as $\sqrt{\mathrm{err}}$, $\Delta\tau$ turns negative in
\NdtauNflip\ of the reparameterisations in Table~\ref{tab:scope}. It also depends on the
range of $T$ the fit covers. We report $\Delta\tau$ and do not rely on it.

\paragraph{Depth and steps (H-D3).}
This is the second registered primary hypothesis. We could not identify a fixed exchange
rate between depth and steps. The registered substitution form assumes that error depends on
depth and steps only through $d\,T^{\kappa}$, as if doubling $T$ multiplied the effective depth
by a fixed factor $2^{\kappa}$. Its
exponent $\kappa$ sits at its upper bound in \NcrsSubFitsAtBound\ of \NcrsSubFits\ fits and
in all \NcrsKappaBootWER\ bootstrap replicates. This form fits worse than the separable one,
by $\Delta\mathrm{AICc} = +\NaiccSSWER$ for WER and $+\NaiccSSSIM$ for SIM-o \citep{aicc}.
The registered criterion was at most $+4$, so
H-D3 fails. Neither form fits the data well. The weighted
$\chi^{2}$ per degree of freedom is \NcrsChiSubWER\ (WER) and \NcrsChiSubSIM\ (SIM-o) for
the substitution form, and \NcrsChiSepWER\ and \NcrsChiSepSIM\ for the separable form. So the
AICc gap compares two misspecified forms.

\paragraph{Shape at fixed parameter count (H-D1).}
Shape matters, but we could not identify how. At $T{=}16$, a fit with separate width and
depth terms beats a fit in the parameter count $N$ alone, with
$\Delta\mathrm{AICc} = \NaiccFNWER$ for WER and $\NaiccFNSIM$ for SIM-o (registered
threshold $-4$). The registered test compares the ratio $\rho = \alpha/\beta$ of the width
and depth exponents across the two metrics. In all four registered fits, and in
\NcrsAmpFitsAtCap\ of \NcrsAmpFits\ fits overall, the coefficient of the width term reaches
the upper limit set for it, so the data cannot pin down the width exponent $\alpha$, and
therefore not $\rho$ either.

\paragraph{Depth optimum and extrapolation (H-D4).}
The WER-best depth is the deepest shape tested in \NcrsDeepCellsWER\ of the
\NcrsDepthCells\ (budget, $T$) cells, so in those cells the grid bounds the WER-optimal depth
only from below. At $T{=}16$, no budget's WER optimum is resolved: the largest margin,
\NcrsDstarMarginCWER\ WER points for the deepest shape of the 125\,M budget, is under two
standard errors of the difference between two configuration means (SE
\NcrsDstarSEdiffCWER), and in the two smaller budgets the two best shapes are within one
standard error of each other. Fits on the two smaller
budgets also do not predict the largest (Fig.~\ref{fig:extrap}). For WER the mean absolute percentage error is
\NmapeWER\%, above the registered 15\%. For SIM-o the fitted form does no better than the
$N$-only model (\NmapeSIM\% against \NmapeNSIM\%).
